\chapter{Advanced Installation}
\label{apendice:instalacao}
% ── index ───────────────────────────────────────────────────────
\index{installation|textbf}
\index{REPL|textbf}
\index{modules}
\index{seed@\texttt{seed()}|textbf}
\index{rnormal@\texttt{rnormal()}|textbf}

% ─────────────────────────────────────────────────────────────────────────────
\section{Compilation with ODBC support}

To connect \hayashi{} to databases via ODBC (SQL Server, Oracle, etc.):

\begin{lstlisting}[style=inline, language=sh]
cargo install hayashi-lang --features odbc
\end{lstlisting}

Requires an ODBC driver installed on the system (\texttt{unixODBC} on Linux).

% ─────────────────────────────────────────────────────────────────────────────
\section{Environment variables}

\begin{center}
\begin{tabular}{ll}
\toprule
\textbf{Variable} & \textbf{Effect} \\
\midrule
\texttt{HAY\_HISTORY}  & Path to REPL history file \\
\texttt{HAY\_NOCOLOR}  & Disables colors in error messages \\
\texttt{HAY\_TRACE}    & Enables execution trace (debugging) \\
\bottomrule
\end{tabular}
\end{center}

% ─────────────────────────────────────────────────────────────────────────────
\section{Modules and optional features}

\subsection*{Advanced econometric models}

The estimators below require the \texttt{greeners} crate as a dependency
of the interpreter. In development builds (\texttt{cargo build}), it is
compiled automatically. No additional installation is required.

\begin{center}
\begin{tabular}{ll}
\toprule
\textbf{Group} & \textbf{Estimators} \\
\midrule
Panel data         & \texttt{fe}, \texttt{re}, \texttt{ab}, \texttt{pcse}, \texttt{gee},
                     \texttt{mundlak}, \texttt{xtgls}, \texttt{mixedlm}, \texttt{clogit} \\
Qualitative response & \texttt{logit}, \texttt{probit}, \texttt{mlogit},
                       \texttt{ologit}, \texttt{oprobit} \\
Count/limited      & \texttt{poisson}, \texttt{nbreg}, \texttt{zinb},
                     \texttt{tobit}, \texttt{heckman} \\
General GLM        & \texttt{glm} (5 families, multiple links) \\
Regularization     & \texttt{lasso}, \texttt{ridge\_reg}, \texttt{enet} \\
Survival           & \texttt{km}, \texttt{cox} \\
Multivariate       & \texttt{pca}, \texttt{factor}, \texttt{cancorr}, \texttt{manova} \\
Time series        & \texttt{var}, \texttt{svar}, \texttt{vecm}, \texttt{garch},
                     \texttt{markov}, \texttt{arima} \\
Filters            & \texttt{hprescott}, \texttt{baxter\_king}, \texttt{holtwinters},
                     \texttt{stl}, \texttt{christiano\_fitzgerald} \\
Causal inference   & \texttt{psm}, \texttt{synth}, \texttt{did}, \texttt{rd} \\
\bottomrule
\end{tabular}
\end{center}

\subsection*{Memory usage}

Computationally intensive models may require stack size adjustment:

\begin{lstlisting}[style=inline, language=sh]
RUST_MIN_STACK=8388608 hay script.hay  // 8 MB of stack
\end{lstlisting}

Relevant for VAR/SVAR with many lags or PCA on large datasets.

% ─────────────────────────────────────────────────────────────────────────────
\section{Editor integration}

\subsection*{Neovim}

Add to \texttt{init.lua} for basic syntax recognition:

\begin{lstlisting}[style=inline, language=Hayashi]
vim.filetype.add({ extension = { hay = "hayashi" } })
\end{lstlisting}

Full highlighting via Tree-sitter is on the roadmap.

\subsection*{VS Code}

Extension available on the marketplace as \texttt{hayashi-lang}.

% ─────────────────────────────────────────────────────────────────────────────
\section{Development builds}

To iterate with local builds (without \texttt{--release}):

\begin{lstlisting}[style=inline, language=sh]
git clone https://github.com/sheep-farm/hayashi
cd hayashi
cargo build            // debug build (recommended for dev)
./target/debug/hay script.hay
\end{lstlisting}

The debug build has extra bounds checks and detailed error messages.
Do not use \texttt{--release} in development — the speed gain does not
compensate for the loss of diagnostic information.

\subsection*{Running scripts}

\begin{lstlisting}[style=inline, language=sh]
./target/debug/hay script.hay        // run script
./target/debug/hay                   // open interactive REPL
./target/debug/hay < script.hay      // read from stdin
\end{lstlisting}
